\section{}

\paragraph{1029年}\eventanchor{LIAO-1029-REGNAL-YEAR}{LIAO-1029-REGNAL-YEAR}　辽以辽圣宗在位第47年作为诏令、赋役与外交文书的纪年框架。账本所列《辽史》〈本纪〉、〈营卫志〉、〈百官志〉；《宋史》辽国相关传；宋人使辽记录为本条来源起点；该条可固定编年位置，具体执行、规模和地方社会影响仍须保留来源差异。

\paragraph{1030年}\eventanchor{LIAO-1030-REGNAL-YEAR}{LIAO-1030-REGNAL-YEAR}　辽以辽圣宗在位第48年作为诏令、赋役与外交文书的纪年框架。账本所列《辽史》〈本纪〉、〈营卫志〉、〈百官志〉；《宋史》辽国相关传；宋人使辽记录为本条来源起点；该条可固定编年位置，具体执行、规模和地方社会影响仍须保留来源差异。

\paragraph{1031年}\eventanchor{LIAO-1031-EVENT-039}{LIAO-1031-EVENT-039}　辽圣宗去世，兴宗即位。账本所列《辽史》〈本纪〉、〈营卫志〉、〈百官志〉；《宋史》辽国相关传；宋人使辽记录为本条来源起点；该条可固定编年位置，具体执行、规模和地方社会影响仍须保留来源差异。

\paragraph{1031年}\eventanchor{LIAO-1031-REGNAL-YEAR}{LIAO-1031-REGNAL-YEAR}　辽以辽圣宗在位第49年作为诏令、赋役与外交文书的纪年框架。账本所列《辽史》〈本纪〉、〈营卫志〉、〈百官志〉；《宋史》辽国相关传；宋人使辽记录为本条来源起点；该条可固定编年位置，具体执行、规模和地方社会影响仍须保留来源差异。

\paragraph{1032年}\eventanchor{LIAO-1032-REGNAL-YEAR}{LIAO-1032-REGNAL-YEAR}　辽以辽兴宗在位第1年作为诏令、赋役与外交文书的纪年框架。账本所列《辽史》〈本纪〉、〈营卫志〉、〈百官志〉；《宋史》辽国相关传；宋人使辽记录为本条来源起点；该条可固定编年位置，具体执行、规模和地方社会影响仍须保留来源差异。

\paragraph{1033年}\eventanchor{LIAO-1033-REGNAL-YEAR}{LIAO-1033-REGNAL-YEAR}　辽以辽兴宗在位第2年作为诏令、赋役与外交文书的纪年框架。账本所列《辽史》〈本纪〉、〈营卫志〉、〈百官志〉；《宋史》辽国相关传；宋人使辽记录为本条来源起点；该条可固定编年位置，具体执行、规模和地方社会影响仍须保留来源差异。

\paragraph{1034年}\eventanchor{LIAO-1034-REGNAL-YEAR}{LIAO-1034-REGNAL-YEAR}　辽以辽兴宗在位第3年作为诏令、赋役与外交文书的纪年框架。账本所列《辽史》〈本纪〉、〈营卫志〉、〈百官志〉；《宋史》辽国相关传；宋人使辽记录为本条来源起点；该条可固定编年位置，具体执行、规模和地方社会影响仍须保留来源差异。

\paragraph{1035年}\eventanchor{LIAO-1035-REGNAL-YEAR}{LIAO-1035-REGNAL-YEAR}　辽以辽兴宗在位第4年作为诏令、赋役与外交文书的纪年框架。账本所列《辽史》〈本纪〉、〈营卫志〉、〈百官志〉；《宋史》辽国相关传；宋人使辽记录为本条来源起点；该条可固定编年位置，具体执行、规模和地方社会影响仍须保留来源差异。

\paragraph{1036年}\eventanchor{LIAO-1036-REGNAL-YEAR}{LIAO-1036-REGNAL-YEAR}　辽以辽兴宗在位第5年作为诏令、赋役与外交文书的纪年框架。账本所列《辽史》〈本纪〉、〈营卫志〉、〈百官志〉；《宋史》辽国相关传；宋人使辽记录为本条来源起点；该条可固定编年位置，具体执行、规模和地方社会影响仍须保留来源差异。

\paragraph{1037年}\eventanchor{LIAO-1037-REGNAL-YEAR}{LIAO-1037-REGNAL-YEAR}　辽以辽兴宗在位第6年作为诏令、赋役与外交文书的纪年框架。账本所列《辽史》〈本纪〉、〈营卫志〉、〈百官志〉；《宋史》辽国相关传；宋人使辽记录为本条来源起点；该条可固定编年位置，具体执行、规模和地方社会影响仍须保留来源差异。

\paragraph{1038年}\eventanchor{LIAO-1038-REGNAL-YEAR}{LIAO-1038-REGNAL-YEAR}　辽以辽兴宗在位第7年作为诏令、赋役与外交文书的纪年框架。账本所列《辽史》〈本纪〉、〈营卫志〉、〈百官志〉；《宋史》辽国相关传；宋人使辽记录为本条来源起点；该条可固定编年位置，具体执行、规模和地方社会影响仍须保留来源差异。

\paragraph{1039年}\eventanchor{LIAO-1039-REGNAL-YEAR}{LIAO-1039-REGNAL-YEAR}　辽以辽兴宗在位第8年作为诏令、赋役与外交文书的纪年框架。账本所列《辽史》〈本纪〉、〈营卫志〉、〈百官志〉；《宋史》辽国相关传；宋人使辽记录为本条来源起点；该条可固定编年位置，具体执行、规模和地方社会影响仍须保留来源差异。

\paragraph{1040年}\eventanchor{LIAO-1040-REGNAL-YEAR}{LIAO-1040-REGNAL-YEAR}　辽以辽兴宗在位第9年作为诏令、赋役与外交文书的纪年框架。账本所列《辽史》〈本纪〉、〈营卫志〉、〈百官志〉；《宋史》辽国相关传；宋人使辽记录为本条来源起点；该条可固定编年位置，具体执行、规模和地方社会影响仍须保留来源差异。

\paragraph{1041年}\eventanchor{LIAO-1041-REGNAL-YEAR}{LIAO-1041-REGNAL-YEAR}　辽以辽兴宗在位第10年作为诏令、赋役与外交文书的纪年框架。账本所列《辽史》〈本纪〉、〈营卫志〉、〈百官志〉；《宋史》辽国相关传；宋人使辽记录为本条来源起点；该条可固定编年位置，具体执行、规模和地方社会影响仍须保留来源差异。

\paragraph{1042年}\eventanchor{LIAO-1042-REGNAL-YEAR}{LIAO-1042-REGNAL-YEAR}　辽以辽兴宗在位第11年作为诏令、赋役与外交文书的纪年框架。账本所列《辽史》〈本纪〉、〈营卫志〉、〈百官志〉；《宋史》辽国相关传；宋人使辽记录为本条来源起点；该条可固定编年位置，具体执行、规模和地方社会影响仍须保留来源差异。

\paragraph{1043年}\eventanchor{LIAO-1043-REGNAL-YEAR}{LIAO-1043-REGNAL-YEAR}　辽以辽兴宗在位第12年作为诏令、赋役与外交文书的纪年框架。账本所列《辽史》〈本纪〉、〈营卫志〉、〈百官志〉；《宋史》辽国相关传；宋人使辽记录为本条来源起点；该条可固定编年位置，具体执行、规模和地方社会影响仍须保留来源差异。

\paragraph{1044年}\eventanchor{LIAO-1044-EVENT-040}{LIAO-1044-EVENT-040}　宋夏议和期间，辽的压力与外交姿态影响三方关系。账本所列《辽史》〈本纪〉、〈营卫志〉、〈百官志〉；《宋史》辽国相关传；宋人使辽记录为本条来源起点；该条可固定编年位置，具体执行、规模和地方社会影响仍须保留来源差异。

\paragraph{1044年}\eventanchor{LIAO-1044-REGNAL-YEAR}{LIAO-1044-REGNAL-YEAR}　辽以辽兴宗在位第13年作为诏令、赋役与外交文书的纪年框架。账本所列《辽史》〈本纪〉、〈营卫志〉、〈百官志〉；《宋史》辽国相关传；宋人使辽记录为本条来源起点；该条可固定编年位置，具体执行、规模和地方社会影响仍须保留来源差异。

\paragraph{1045年}\eventanchor{LIAO-1045-REGNAL-YEAR}{LIAO-1045-REGNAL-YEAR}　辽以辽兴宗在位第14年作为诏令、赋役与外交文书的纪年框架。账本所列《辽史》〈本纪〉、〈营卫志〉、〈百官志〉；《宋史》辽国相关传；宋人使辽记录为本条来源起点；该条可固定编年位置，具体执行、规模和地方社会影响仍须保留来源差异。

\paragraph{1046年}\eventanchor{LIAO-1046-REGNAL-YEAR}{LIAO-1046-REGNAL-YEAR}　辽以辽兴宗在位第15年作为诏令、赋役与外交文书的纪年框架。账本所列《辽史》〈本纪〉、〈营卫志〉、〈百官志〉；《宋史》辽国相关传；宋人使辽记录为本条来源起点；该条可固定编年位置，具体执行、规模和地方社会影响仍须保留来源差异。

\paragraph{1047年}\eventanchor{LIAO-1047-REGNAL-YEAR}{LIAO-1047-REGNAL-YEAR}　辽以辽兴宗在位第16年作为诏令、赋役与外交文书的纪年框架。账本所列《辽史》〈本纪〉、〈营卫志〉、〈百官志〉；《宋史》辽国相关传；宋人使辽记录为本条来源起点；该条可固定编年位置，具体执行、规模和地方社会影响仍须保留来源差异。

\paragraph{1048年}\eventanchor{LIAO-1048-REGNAL-YEAR}{LIAO-1048-REGNAL-YEAR}　辽以辽兴宗在位第17年作为诏令、赋役与外交文书的纪年框架。账本所列《辽史》〈本纪〉、〈营卫志〉、〈百官志〉；《宋史》辽国相关传；宋人使辽记录为本条来源起点；该条可固定编年位置，具体执行、规模和地方社会影响仍须保留来源差异。

\paragraph{1049年}\eventanchor{LIAO-1049-REGNAL-YEAR}{LIAO-1049-REGNAL-YEAR}　辽以辽兴宗在位第18年作为诏令、赋役与外交文书的纪年框架。账本所列《辽史》〈本纪〉、〈营卫志〉、〈百官志〉；《宋史》辽国相关传；宋人使辽记录为本条来源起点；该条可固定编年位置，具体执行、规模和地方社会影响仍须保留来源差异。

\paragraph{1050年}\eventanchor{LIAO-1050-REGNAL-YEAR}{LIAO-1050-REGNAL-YEAR}　辽以辽兴宗在位第19年作为诏令、赋役与外交文书的纪年框架。账本所列《辽史》〈本纪〉、〈营卫志〉、〈百官志〉；《宋史》辽国相关传；宋人使辽记录为本条来源起点；该条可固定编年位置，具体执行、规模和地方社会影响仍须保留来源差异。

\paragraph{1051年}\eventanchor{LIAO-1051-REGNAL-YEAR}{LIAO-1051-REGNAL-YEAR}　辽以辽兴宗在位第20年作为诏令、赋役与外交文书的纪年框架。账本所列《辽史》〈本纪〉、〈营卫志〉、〈百官志〉；《宋史》辽国相关传；宋人使辽记录为本条来源起点；该条可固定编年位置，具体执行、规模和地方社会影响仍须保留来源差异。

\paragraph{1052年}\eventanchor{LIAO-1052-REGNAL-YEAR}{LIAO-1052-REGNAL-YEAR}　辽以辽兴宗在位第21年作为诏令、赋役与外交文书的纪年框架。账本所列《辽史》〈本纪〉、〈营卫志〉、〈百官志〉；《宋史》辽国相关传；宋人使辽记录为本条来源起点；该条可固定编年位置，具体执行、规模和地方社会影响仍须保留来源差异。

\paragraph{1053年}\eventanchor{LIAO-1053-REGNAL-YEAR}{LIAO-1053-REGNAL-YEAR}　辽以辽兴宗在位第22年作为诏令、赋役与外交文书的纪年框架。账本所列《辽史》〈本纪〉、〈营卫志〉、〈百官志〉；《宋史》辽国相关传；宋人使辽记录为本条来源起点；该条可固定编年位置，具体执行、规模和地方社会影响仍须保留来源差异。

\paragraph{1054年}\eventanchor{LIAO-1054-REGNAL-YEAR}{LIAO-1054-REGNAL-YEAR}　辽以辽兴宗在位第23年作为诏令、赋役与外交文书的纪年框架。账本所列《辽史》〈本纪〉、〈营卫志〉、〈百官志〉；《宋史》辽国相关传；宋人使辽记录为本条来源起点；该条可固定编年位置，具体执行、规模和地方社会影响仍须保留来源差异。

\paragraph{1055年}\eventanchor{LIAO-1055-EVENT-041}{LIAO-1055-EVENT-041}　辽兴宗去世，道宗即位。账本所列《辽史》〈本纪〉、〈营卫志〉、〈百官志〉；《宋史》辽国相关传；宋人使辽记录为本条来源起点；该条可固定编年位置，具体执行、规模和地方社会影响仍须保留来源差异。

\paragraph{1055年}\eventanchor{LIAO-1055-REGNAL-YEAR}{LIAO-1055-REGNAL-YEAR}　辽以辽兴宗在位第24年作为诏令、赋役与外交文书的纪年框架。账本所列《辽史》〈本纪〉、〈营卫志〉、〈百官志〉；《宋史》辽国相关传；宋人使辽记录为本条来源起点；该条可固定编年位置，具体执行、规模和地方社会影响仍须保留来源差异。

\paragraph{1056年}\eventanchor{LIAO-1056-REGNAL-YEAR}{LIAO-1056-REGNAL-YEAR}　辽以辽道宗在位第1年作为诏令、赋役与外交文书的纪年框架。账本所列《辽史》〈本纪〉、〈营卫志〉、〈百官志〉；《宋史》辽国相关传；宋人使辽记录为本条来源起点；该条可固定编年位置，具体执行、规模和地方社会影响仍须保留来源差异。

\paragraph{1057年}\eventanchor{LIAO-1057-REGNAL-YEAR}{LIAO-1057-REGNAL-YEAR}　辽以辽道宗在位第2年作为诏令、赋役与外交文书的纪年框架。账本所列《辽史》〈本纪〉、〈营卫志〉、〈百官志〉；《宋史》辽国相关传；宋人使辽记录为本条来源起点；该条可固定编年位置，具体执行、规模和地方社会影响仍须保留来源差异。

\paragraph{1058年}\eventanchor{LIAO-1058-REGNAL-YEAR}{LIAO-1058-REGNAL-YEAR}　辽以辽道宗在位第3年作为诏令、赋役与外交文书的纪年框架。账本所列《辽史》〈本纪〉、〈营卫志〉、〈百官志〉；《宋史》辽国相关传；宋人使辽记录为本条来源起点；该条可固定编年位置，具体执行、规模和地方社会影响仍须保留来源差异。

\paragraph{1059年}\eventanchor{LIAO-1059-REGNAL-YEAR}{LIAO-1059-REGNAL-YEAR}　辽以辽道宗在位第4年作为诏令、赋役与外交文书的纪年框架。账本所列《辽史》〈本纪〉、〈营卫志〉、〈百官志〉；《宋史》辽国相关传；宋人使辽记录为本条来源起点；该条可固定编年位置，具体执行、规模和地方社会影响仍须保留来源差异。

\paragraph{1060年}\eventanchor{LIAO-1060-REGNAL-YEAR}{LIAO-1060-REGNAL-YEAR}　辽以辽道宗在位第5年作为诏令、赋役与外交文书的纪年框架。账本所列《辽史》〈本纪〉、〈营卫志〉、〈百官志〉；《宋史》辽国相关传；宋人使辽记录为本条来源起点；该条可固定编年位置，具体执行、规模和地方社会影响仍须保留来源差异。

\paragraph{1061年}\eventanchor{LIAO-1061-REGNAL-YEAR}{LIAO-1061-REGNAL-YEAR}　辽以辽道宗在位第6年作为诏令、赋役与外交文书的纪年框架。账本所列《辽史》〈本纪〉、〈营卫志〉、〈百官志〉；《宋史》辽国相关传；宋人使辽记录为本条来源起点；该条可固定编年位置，具体执行、规模和地方社会影响仍须保留来源差异。

\paragraph{1062年}\eventanchor{LIAO-1062-REGNAL-YEAR}{LIAO-1062-REGNAL-YEAR}　辽以辽道宗在位第7年作为诏令、赋役与外交文书的纪年框架。账本所列《辽史》〈本纪〉、〈营卫志〉、〈百官志〉；《宋史》辽国相关传；宋人使辽记录为本条来源起点；该条可固定编年位置，具体执行、规模和地方社会影响仍须保留来源差异。

\paragraph{1063年}\eventanchor{LIAO-1063-REGNAL-YEAR}{LIAO-1063-REGNAL-YEAR}　辽以辽道宗在位第8年作为诏令、赋役与外交文书的纪年框架。账本所列《辽史》〈本纪〉、〈营卫志〉、〈百官志〉；《宋史》辽国相关传；宋人使辽记录为本条来源起点；该条可固定编年位置，具体执行、规模和地方社会影响仍须保留来源差异。

\paragraph{1064年}\eventanchor{LIAO-1064-REGNAL-YEAR}{LIAO-1064-REGNAL-YEAR}　辽以辽道宗在位第9年作为诏令、赋役与外交文书的纪年框架。账本所列《辽史》〈本纪〉、〈营卫志〉、〈百官志〉；《宋史》辽国相关传；宋人使辽记录为本条来源起点；该条可固定编年位置，具体执行、规模和地方社会影响仍须保留来源差异。

\paragraph{1065年}\eventanchor{LIAO-1065-REGNAL-YEAR}{LIAO-1065-REGNAL-YEAR}　辽以辽道宗在位第10年作为诏令、赋役与外交文书的纪年框架。账本所列《辽史》〈本纪〉、〈营卫志〉、〈百官志〉；《宋史》辽国相关传；宋人使辽记录为本条来源起点；该条可固定编年位置，具体执行、规模和地方社会影响仍须保留来源差异。

\paragraph{1066年}\eventanchor{LIAO-1066-REGNAL-YEAR}{LIAO-1066-REGNAL-YEAR}　辽以辽道宗在位第11年作为诏令、赋役与外交文书的纪年框架。账本所列《辽史》〈本纪〉、〈营卫志〉、〈百官志〉；《宋史》辽国相关传；宋人使辽记录为本条来源起点；该条可固定编年位置，具体执行、规模和地方社会影响仍须保留来源差异。

\paragraph{1067年}\eventanchor{LIAO-1067-REGNAL-YEAR}{LIAO-1067-REGNAL-YEAR}　辽以辽道宗在位第12年作为诏令、赋役与外交文书的纪年框架。账本所列《辽史》〈本纪〉、〈营卫志〉、〈百官志〉；《宋史》辽国相关传；宋人使辽记录为本条来源起点；该条可固定编年位置，具体执行、规模和地方社会影响仍须保留来源差异。

\paragraph{1068年}\eventanchor{LIAO-1068-REGNAL-YEAR}{LIAO-1068-REGNAL-YEAR}　辽以辽道宗在位第13年作为诏令、赋役与外交文书的纪年框架。账本所列《辽史》〈本纪〉、〈营卫志〉、〈百官志〉；《宋史》辽国相关传；宋人使辽记录为本条来源起点；该条可固定编年位置，具体执行、规模和地方社会影响仍须保留来源差异。

\paragraph{1069年}\eventanchor{LIAO-1069-REGNAL-YEAR}{LIAO-1069-REGNAL-YEAR}　辽以辽道宗在位第14年作为诏令、赋役与外交文书的纪年框架。账本所列《辽史》〈本纪〉、〈营卫志〉、〈百官志〉；《宋史》辽国相关传；宋人使辽记录为本条来源起点；该条可固定编年位置，具体执行、规模和地方社会影响仍须保留来源差异。

\paragraph{1070年}\eventanchor{LIAO-1070-REGNAL-YEAR}{LIAO-1070-REGNAL-YEAR}　辽以辽道宗在位第15年作为诏令、赋役与外交文书的纪年框架。账本所列《辽史》〈本纪〉、〈营卫志〉、〈百官志〉；《宋史》辽国相关传；宋人使辽记录为本条来源起点；该条可固定编年位置，具体执行、规模和地方社会影响仍须保留来源差异。

\paragraph{1071年}\eventanchor{LIAO-1071-REGNAL-YEAR}{LIAO-1071-REGNAL-YEAR}　辽以辽道宗在位第16年作为诏令、赋役与外交文书的纪年框架。账本所列《辽史》〈本纪〉、〈营卫志〉、〈百官志〉；《宋史》辽国相关传；宋人使辽记录为本条来源起点；该条可固定编年位置，具体执行、规模和地方社会影响仍须保留来源差异。

\paragraph{1072年}\eventanchor{LIAO-1072-REGNAL-YEAR}{LIAO-1072-REGNAL-YEAR}　辽以辽道宗在位第17年作为诏令、赋役与外交文书的纪年框架。账本所列《辽史》〈本纪〉、〈营卫志〉、〈百官志〉；《宋史》辽国相关传；宋人使辽记录为本条来源起点；该条可固定编年位置，具体执行、规模和地方社会影响仍须保留来源差异。

\paragraph{1073年}\eventanchor{LIAO-1073-REGNAL-YEAR}{LIAO-1073-REGNAL-YEAR}　辽以辽道宗在位第18年作为诏令、赋役与外交文书的纪年框架。账本所列《辽史》〈本纪〉、〈营卫志〉、〈百官志〉；《宋史》辽国相关传；宋人使辽记录为本条来源起点；该条可固定编年位置，具体执行、规模和地方社会影响仍须保留来源差异。

\paragraph{1074年}\eventanchor{LIAO-1074-REGNAL-YEAR}{LIAO-1074-REGNAL-YEAR}　辽以辽道宗在位第19年作为诏令、赋役与外交文书的纪年框架。账本所列《辽史》〈本纪〉、〈营卫志〉、〈百官志〉；《宋史》辽国相关传；宋人使辽记录为本条来源起点；该条可固定编年位置，具体执行、规模和地方社会影响仍须保留来源差异。

\paragraph{1075年}\eventanchor{LIAO-1075-REGNAL-YEAR}{LIAO-1075-REGNAL-YEAR}　辽以辽道宗在位第20年作为诏令、赋役与外交文书的纪年框架。账本所列《辽史》〈本纪〉、〈营卫志〉、〈百官志〉；《宋史》辽国相关传；宋人使辽记录为本条来源起点；该条可固定编年位置，具体执行、规模和地方社会影响仍须保留来源差异。

